\DocumentMetadata{
  pdfversion=2.0,
  pdfstandard=ua-2,
  lang=en-US,
  tagging=on,
  tagging-setup={math/setup=mathml-SE,tabsorder=structure}
}
\documentclass[11pt,letterpaper]{article}
\usepackage[margin=0.68in,includefoot]{geometry}
\usepackage{newcomputermodern}
\usepackage{amsmath,amscd,mathtools}
\usepackage{tikz,adjustbox}
\providecommand{\square}{\mdlgwhtsquare}
\usepackage[hidelinks]{hyperref}
\hypersetup{
  pdftitle={Combinatorics PhD Exam, May 1991},
  pdfauthor={University of Florida Department of Mathematics},
  pdfsubject={Graduate mathematics examination},
  pdfdisplaydoctitle=true,
  bookmarksopen=true
}
\setcounter{secnumdepth}{0}
\setlength{\parindent}{0pt}
\setlength{\parskip}{0.28em}
\setlength{\emergencystretch}{3em}
\setlength{\leftmargini}{2.15em}
\setlength{\leftmarginii}{2.65em}
\setlength{\leftmarginiii}{3em}
\pagestyle{empty}
\raggedbottom
\allowdisplaybreaks
\NewTaggingSocketPlug{block/list/label}{examproblem}{%
  \tagstructbegin{tag=Lbl}\tagstructend%
  \tagstructbegin{tag=\LBody}%
  \tagstructbegin{tag=\ProblemHeadingTag,title-o={\ProblemHeadingText}}%
  \tagmcbegin{tag=\ProblemHeadingTag,actualtext={\ProblemHeadingText}}%
  #2%
  \tagmcend%
  \tagstructend%
}
\newenvironment{examproblems}[2]{%
  \def\ProblemHeadingTag{#2}%
  \AssignTaggingSocketPlug{block/list/label}{examproblem}%
  \begin{enumerate}%
  \setcounter{enumi}{#1}%
  \renewcommand{\labelenumi}{\textbf{\theenumi.}}%
  \setlength{\topsep}{0.35em}%
  \setlength{\partopsep}{0pt}%
  \setlength{\itemsep}{0.48em}%
  \setlength{\parsep}{0.18em}%
}{\end{enumerate}}
\newenvironment{examsubparts}{%
  \AssignTaggingSocketPlug{block/list/label}{default}%
  \begin{enumerate}%
  \setlength{\topsep}{0.18em}%
  \setlength{\partopsep}{0pt}%
  \setlength{\itemsep}{0.16em}%
  \setlength{\parsep}{0.1em}%
}{\end{enumerate}}
\newenvironment{examinstructions}{%
  \par\begingroup\itshape
}{%
  \par\endgroup\vspace{0.18em}
}
\makeatletter
\renewcommand\subsection{\@startsection{subsection}{2}{0pt}%
  {-1.25ex plus -.25ex minus -.1ex}%
  {0.55ex plus .1ex}%
  {\normalfont\large\bfseries}}
\renewcommand{\@maketitle}{%
  \newpage\null\vskip -1em%
  {\footnotesize\bfseries
    \href{https://www.ufl.edu/}{University of Florida}, \href{https://math.ufl.edu/}{Department of Mathematics}\par}%
  \vskip -0.5em%
  \tagstructbegin{tag=H1,title={Combinatorics PhD Exam, May 1991}}%
  \tagpdfparaOff%
  \tagmcbegin{tag=H1}%
    {\LARGE\bfseries\href{https://gma.math.ufl.edu/exams/combinatorics-phd/}{Combinatorics PhD Exam}, May 1991\par}%
  \tagmcend%
  \tagpdfparaOn%
  \tagstructend%
  \vskip 0.25em%
  \tagmcbegin{artifact}%
    \hrule height 1pt%
  \tagmcend%
  \vskip 1em%
}
\makeatother
\title{Combinatorics PhD Exam, May 1991}
\author{}
\date{}
\begin{document}
\maketitle
\thispagestyle{empty}
\begin{examinstructions}
Do 9 out of 12 problems. Show all of your work.
\end{examinstructions}
\begin{examproblems}{0}{H2}
\def\ProblemHeadingText{Problem 1.}
\item
How many~\(k\)-tuples \((A_1,A_2,\ldots,A_k)\) of nested subsets of an~\(n\)-set are there, where \(A_1\subseteq A_2\subseteq\cdots\subseteq A_k\)?
\def\ProblemHeadingText{Problem 2.}
\item
How many words of length~\(n\) are there in the alphabet \((a,b,c)\) with no two adjacent~\(a\)’s? With no two adjacent letters equal?
\def\ProblemHeadingText{Problem 3.}
\item
Find the generating function for the sequence \(\{a_n\}\) given by the recurrence
\[
a_n=a_{n-1}+6a_{n-2}+2^n
\]
 where \(a_0=1\) and \(a_1=3\).
\def\ProblemHeadingText{Problem 4.}
\item
Bipartite graphs:
\begin{examsubparts}
\item[\textnormal{(a)}]
Determine for which values of~\(m\) and~\(n\) the complete bipartite graph \(K_{mn}\) is
\begin{examsubparts}
\item[\textnormal{(1)}]
planar;
\item[\textnormal{(2)}]
Eulerian;
\item[\textnormal{(3)}]
Hamiltonian.
\end{examsubparts}
\item[\textnormal{(b)}]
Prove: A plane graph where every face has an even number of edges must be bipartite.
\end{examsubparts}
\def\ProblemHeadingText{Problem 5.}
\item
Prove that
\[
\sum_{k=1}^{n} S(n,k)x(x-1)\cdots(x-k+1)=x^n.
\]
\def\ProblemHeadingText{Problem 6.}
\item
Use Möbius inversion to find the number of integers between~\(1\) and~\(n\) and coprime to~\(n\), given the prime factorization of~\(n\).
\def\ProblemHeadingText{Problem 7.}
\item
Prove or disprove: There is a matching from~\(V\) to~\(W\) in a bipartite graph if, for some fixed~\(k\), there are~\(k\) or more edges incident with each vertex of~\(V\) and~\(k\) or fewer edges incident with each vertex of~\(W\).
\def\ProblemHeadingText{Problem 8.}
\item
Let \(A_{m,n}\) be the collection of all sets \(X=\{N_1,N_2,\ldots,N_k\}\), where~\(k\) is a non-negative integer, and for each \(i\leq k\), \(N_i\) is a positive integer divisor of \(N=2^m3^n\), and where \(N_i\mid N_j\) only if \(i=j\). Define a partial order on \(A_{m,n}\), where \(X=\{N_1,N_2,\ldots,N_k\}\), \(X'=\{N'_1,N'_2,\ldots,N'_l\}\), by \(X\preceq X'\) if and only if \(N_i\in X\Rightarrow\) there exists \(N'_j\in X'\) with \(N_i\mid N'_j\). Prove that \(A_{m,n}\) is a distributive lattice, and determine its poset of join-irreducible elements.
\def\ProblemHeadingText{Problem 9.}
\item
Let~\(I\) be a collection of subsets of~\(E\), where~\(E\) is finite, and suppose that \(I\neq\varnothing\) and that \(Y\in I,\ X\subseteq Y\Rightarrow X\in I\). Let \(c:E\to\mathbb{R}^+\) assign a positive real number to each element of~\(E\). Consider the optimization problem:

\par
Find \(\max_{X\in I}\sum_{x\in X}c(x)\)

\par
Prove that the greedy algorithm solves this problem if and only if~\(I\) is the collection of independent sets of some matroid on~\(E\).
\def\ProblemHeadingText{Problem 10.}
\item
State the MacWilliams identity, and use it to compute the weight polynomial of the \((15,11)\) Hamming binary code.
\def\ProblemHeadingText{Problem 11.}
\item
A \((v,k,\lambda)\)-design~\(\Sigma\) is an incidence structure consisting of a set~\(S\) of~\(v\) points and a collection of~\(k\)-subsets of~\(S\) called blocks, such that every pair of points of~\(S\) is contained in precisely~\(\lambda\) blocks.

\par
Prove that, in the range \(30\leq v\leq45\), a \((v,6,1)\)-design exists if and only if \(v=31\).
\def\ProblemHeadingText{Problem 12.}
\item
Use the Hall Multiplier Theorem
\begin{examsubparts}
\item[\textnormal{(i)}]
to prove the non-existence of a cyclic \((31,10,3)\) difference set.
\item[\textnormal{(ii)}]
to find a cyclic \((31,6,1)\) difference set.
\end{examsubparts}
\end{examproblems}
\end{document}
